\par\Needspace{7\baselineskip}\bigskip
\begin{center}
\pdfbookmark[2]{Olbers an Gauss, Bremen, 5. April 1812}{pallas-brief-18120405}
\textsc{Olbers} an \textsc{Gauss}, Bremen, 5. April 1812.
\end{center}
\noindent\hspace*{\parindent}\{. . . . .\quad Ganz entzückt bin ich über die äusserst merkwürdigen Resultate, die Sie mir aus Ihren tiefsinnigen Untersuchungen über die Pallasbahn mittheilen. Allerdings halte auch ich Ihre Entdeckung, dass \(18\text{\AstroSymbols ♃}\) genau \(=7\text{\AstroSymbols ⚴}\) sei, für eine der merkwürdigsten, die seit langer Zeit im Sonnensystem gemacht worden ist. — Aber wie ist es mit der Ceres? Hat diese eine beträchtliche von \(18\text{\AstroSymbols ♃}-7\text{\AstroSymbols ⚳}\) abhängende Gleichung? Wird sich für diese und die übrigen kleinen Planeten nicht auch ähnliche Gleichungen von der Form \(m\text{\AstroSymbols ♃}=n\text{\AstroSymbols ⚳}\) u. s. w. finden lassen?\quad . . . . .\}
\par\Needspace{7\baselineskip}\bigskip
\begin{center}
\pdfbookmark[2]{Gauss an Bessel, Göttingen, 5. Mai 1812}{pallas-brief-18120505}
\textsc{Gauss} an \textsc{Bessel}, Göttingen, 5. Mai 1812.
\end{center}
\noindent\hspace*{\parindent}. . . . .\ habe ich mich hauptsächlich mit den Pallasstörungen durch Jupiter beschäftigt. Sie werden darüber in Nr. 67 unserer Gelehrten Anzeigen einiges gelesen haben. Ihnen theile ich das merkwürdige daselbst in einer Chiffre niedergelegte Resultat gern mit, doch mit der Bitte, dass es vorerst \textls[180]{ganz unter uns beiden} bleibe. Es besteht darin, dass die mittlern Bewegungen von {\AstroSymbols ♃} und {\AstroSymbols ⚴} in dem \textls[180]{rationalen} Verhältniss \(7:18\) stehen, was sich durch die Einwirkung Jupiters immer \textls[180]{genau} wieder herstellt, wie die Rotationszeit unsers Mondes. Ich habe mit einer zweiten Berechnung der periodischen Störungen bereits einen Anfang gemacht.\quad . . . . .
\par\Needspace{7\baselineskip}\bigskip
\begin{center}
\pdfbookmark[2]{Olrers an Gauss, Paris, 12. Mai 1812}{pallas-brief-18120512}
\textsc{Olrers} an \textsc{Gauss}, Paris, 12. Mai 1812.
\end{center}
\noindent\hspace*{\parindent}\{. . . . .\quad \textsc{Laplace} erkundigt sich sehr umständlich nach Ihnen und Ihrer Arbeit über die Pallas. Er scheint sehr neugierig auf Ihre Methode zu sein, von der ich ihm nichts sagen konnte, auch nichts sagen wollte. Sie können nicht glauben, wie enthusiastisch er und alle Andern . . . . .\ die Hoffnung aufnahmen, Sie . . . . .\ bei einer meiner künftigen Reisen nach Paris zur Begleitung bereden zu können.\quad . . . . .\}
\par
